\documentclass[10pt,a4paper]{article}
\usepackage{amsmath,amssymb,amsthm,graphicx}
\usepackage{hyperref}
\usepackage{latexsym}
\usepackage{amssymb}
\usepackage{amsthm}
\usepackage{ulem}
\usepackage{enumitem}
\setlist{itemsep=0.4pt,topsep=0pt}
\parindent 0pt
\parindent 0pt
\parskip 12pt
\newcommand{\HCF}{\mathrm{hcf}}
\newcommand{\ZZ}{\mathbb{Z}}
\newcommand{\rk}{\mathrm{rk}}
\DeclareMathOperator{\len}{length}
\DeclareMathOperator{\bl}{b}
\DeclareMathOperator{\cl}{cl}
\DeclareMathOperator{\di}{d}
\DeclareMathOperator{\Cl}{Cl}
\DeclareMathOperator{\acl}{acl}
\DeclareMathOperator{\tp}{tp}
\DeclareMathOperator{\Th}{Th}
\newcommand{\Cor}{\mathcal{C}_\mathrm{or}}
\newcommand{\Cun}{\mathcal{C}_\mathrm{un}}
\newcommand{\mons}{\mathbb{M}}
\newcommand{\strong}{\mathrm{strong}}
\newcommand{\axioms}{\Delta_{\mathrm{FSTS}}}
\newtheoremstyle{note}% hnamei
{8pt}% hSpace abovei
{14pt}% hSpace belowi
{\normalfont}% hBody fonti
{}% hIndent amounti1
{\bfseries}% hTheorem head fonti
{ }% hPunctuation after theorem headi
{.5em}% hSpace after theorem headi2
{}% hTheorem head spec (can be left empty, meaning ‘normal’)i
\theoremstyle{note}
\newtheorem{thm}{Theorem}[section]
%\newtheorem{theorem}{Theorem}
\newtheorem{lemma}[thm]{Lemma} 
\newtheorem{df}[thm]{Definition}
\newtheorem{rmk}[thm]{Remark} 
\newtheorem{prop}[thm]{Proposition}
\newtheorem{fact}[thm]{Fact}
\newtheorem{cor}[thm]{Corollary}
\newtheorem{conj}[thm]{Conjecture}
\newtheorem{ex}[thm]{Example}
\newtheorem{qu}[thm]{Question}
\newtheorem*{notation}{Notation} 
\title{The binary tree carries an HF-ordering whose smallest element is the root}
\author{Silvia Barbina, Enrique Casanovas}
\date{arxiv:2411.13723v3; Creative Commons Attribution 4.0 International}
\begin{document}
\maketitle
\noindent\textit{Editorial note.} The counted statement is the author's labelled environment at source0.tex lines 820-821 and its proof environment opens at line 822; the pairing is by adjacency and the retained proof range is the whole environment.\par\medskip
\noindent\textit{Editorial note.} This source is newly approved for this line; the preamble is built from the author's own package lines, macro block and theorem-environment declarations (source lines 4-65), with the author's optioned amsmath line omitted so that it cannot clash with the wrapper's amsmath.\par\medskip
\noindent\textit{Editorial note.} The source carries no label-number dump and the retained counter values are left empty, so the printed numbering of the card is the wrapper's default; the identity of the counted item is its author label.\par\medskip
\noindent\textit{Editorial note.} Metadata provenance: the source record of the pinned copy and the licence evidence file are both present; the licence is CC BY 4.0 with status supported.\par\medskip
\noindent\textit{Editorial note.} r002 note: the pre-publication gate compares the card's source\_version with the source\_version recorded in the licence evidence of the same source. This source's evidence file records exactly "2411.13723v3", while the card had carried "arXiv:2411.13723v3 (CC BY 4.0)" copied from the source record; r002 aligns the card with the evidence and changes nothing else. The evidence file itself is not modified.\par\medskip
\noindent\textit{Editorial scope.} One original result of the pinned author manuscript, counted once: an HF-ordering of the binary tree with the root as least element (author label April\_11\_2024\_2).\par\medskip
\noindent\textit{Allowed standard prerequisites (editorial).}\begin{itemize}
\item Binary tree of a given height. The binary tree of height alpha is the set of finite sequences of zeros and ones of length below alpha, ordered by extension.. It is the structure that is ordered.
\item Natural HF-ordering. The natural ordering of the tree puts the root at the top, so that every node is below the root and the blocks are formed by a node and its two immediate successors.. It is the ordering that the proof reverses.
\item HF-ordering. An HF-ordering is a linear ordering in which each element has at most one pair of smaller elements forming a block with it.. It is the property that has to survive the reversal.
\end{itemize}\medskip
\begin{lemma}\label{April_11_2024_2} The binary tree of height $\alpha$ has an HF-ordering where the root  $a_\emptyset$ is the smallest element instead of being the largest.
\end{lemma}

\begin{proof}  If $<$ is  a natural HF-ordering of the binary tree, with $a_\emptyset$ at the top, then the reverse order $>$ is also an HF-ordering and now $a_\emptyset$ is the smallest element.
\end{proof}

\end{document}
